\documentclass[letterpaper]{article}
\usepackage{amsmath,amssymb}
\usepackage[letterpaper,margin=1in]{geometry}
\setlength{\parindent}{0pt}

\begin{document}

\textbf{MATH 1564 --- Notes 09/22}

\bigskip

\textbf{Vector spaces}

\medskip

Vectors are things we can add and scale.

\medskip

\textbf{Examples.}
\begin{itemize}
\item The set of $m\times n$ matrices: if $A,B\in M_{m\times n}$, then $2A+B\in M_{m\times n}$.
\item Polynomials. If
\[
f(x)=a_0+a_1x+\cdots+a_nx^n,
\qquad
g(x)=b_0+b_1x+\cdots+b_mx^m,
\qquad m<n,
\]
then
\[
(f+g)(x)=a_0+b_0+(a_1+b_1)x+\cdots+(a_m+b_m)x^m+a_{m+1}x^{m+1}+\cdots+a_nx^n,
\]
and
\[
2f(x)=2a_0+\cdots+2a_nx^n.
\]
\item Sequences of numbers $(x_n)$, where $n$ is an integer. For example,
\[
2^n=\{\ldots,2^{-1},2^0,2^1,2^2,\ldots\},
\]
and addition and scalar multiplication are defined term-by-term:
\[
(2x+y)_n=2x_n+y_n.
\]
\item Functions $f:X\to\mathbb{R}$, with
\[
(f+g)(x)=f(x)+g(x),
\qquad
(2f)(x)=2f(x).
\]
\end{itemize}

\newpage

\textbf{Definition.} Let $V$ be a set whose elements are called vectors, with two operations:
\begin{enumerate}
\item $V\times V\to V$, denoted $(v,w)\mapsto v+w$ \textbf{(addition)};
\item $\mathbb{R}\times V\to V$, denoted $(c,v)\mapsto cv$ \textbf{(scaling)}.
\end{enumerate}

\medskip

For all $u,v,w\in V$ and $c,d\in\mathbb{R}$, the following vector-space properties hold:
\begin{enumerate}
\item $u+v=v+u$.
\item $(u+v)+w=u+(v+w)$.
\item There is a zero vector $0$ with $v+0=v$ for every $v\in V$.
\item For every $v\in V$, there is $-v$ with $v+(-v)=0$.
\item $c(u+v)=cu+cv$.
\item $(c+d)u=cu+du$.
\item $(cd)u=c(du)$.
\item $1v=v$.
\end{enumerate}

\medskip

Useful consequences:
\begin{enumerate}
\setcounter{enumi}{8}
\item The zero vector is unique.
\item $v$ uniquely determines $-v$.
\item $0v=0$.
\item $c0=0$.
\item $-v=(-1)v$.
\end{enumerate}

\medskip

\textbf{Proof of (9).} If $0'$ is another zero vector, then
\[
0=0+0'=0'.
\]

\textbf{Proof of (10).} If $-v$ and $v'$ both satisfy $-v+v=0=v'+v$, then
\[
-v=-v+0=-v+(v+v')=0+v'=v'.
\]

\newpage

\textbf{Example.} Any subspace of $\mathbb{R}^n$ is a vector space.

\medskip

\textbf{Note.} In the definition of a vector space, the scalars need not come from $\mathbb{R}$. One can use another field, such as $\mathbb{C}$ or $\mathbb{Q}$.

\medskip

There is also a related notion of a \textbf{module}, where scalars can come from objects such as $\mathbb{Z}$ or $\mathbb{Z}/n\mathbb{Z}$ when $n$ is not prime, or from polynomials.

\bigskip

\textbf{Topics in vector spaces}
\begin{enumerate}
\item Linear combinations and subspaces.
\item Linear independence and span.
\item Linear maps, the image of $T$, and the kernel of $T$.
\item Bases, coordinates, and dimension.
\end{enumerate}

\newpage

\textbf{Definition.} Let $S$ be a set of vectors in a vector space $V$. The \textbf{span} of $S$ is the set of all linear combinations of vectors in $S$.

\medskip

The set $S$ need not be finite. A linear combination has the form
\[
c_1v_1+\cdots+c_kv_k.
\]

\medskip

\textbf{Example.}
\[
S=\{1,x^2,x^4,\ldots\}.
\]
Then $\operatorname{span}S$ is the set of polynomials with no odd monomials.

\bigskip

\textbf{Definition.} A \textbf{basis} of a vector space $V$ is a possibly infinite set $\beta$ of vectors in $V$ such that
\[
V=\operatorname{span}\beta
\]
and, for every finite selection $v_1,\ldots,v_k\in\beta$,
\[
c_1v_1+\cdots+c_kv_k=0
\quad\Longrightarrow\quad
\begin{bmatrix}c_1\\ \vdots\\ c_k\end{bmatrix}
=
\begin{bmatrix}0\\ \vdots\\ 0\end{bmatrix}.
\]

\medskip

\textbf{Example.} A basis for the vector space of all polynomials with real coefficients is
\[
\beta=\{1,x,x^2,\ldots,x^n,\ldots\}.
\]

\textbf{Theorem.} The set $\{1,x,x^2,\ldots,x^n,\ldots\}$ is independent.

\textbf{Proof.} Suppose
\[
c_1x^{i_1}+\cdots+c_kx^{i_k}=0,
\qquad i_1<i_2<\cdots<i_k.
\]
Divide by $x^{i_1}$ to obtain
\[
c_1+c_2x^{i_2-i_1}+\cdots+c_kx^{i_k-i_1}=0.
\]
Setting $x=0$ gives $c_1=0$; induction on $k$ finishes the proof.

\newpage

\textbf{Corollary.} The set
\[
\{1,x,x^2,\ldots,x^n,\ldots\}
\]
is a basis of $V=\mathbb{R}[x]$, the vector space of polynomials with real coefficients.

\bigskip

\textbf{Definition.} A \textbf{subspace} of a vector space is a subset $S$ such that
\[
x,y\in S\Longrightarrow x+y\in S,
\qquad
c\in\mathbb{R},\ x\in S\Longrightarrow cx\in S.
\]

\textbf{Note.} A subspace of a vector space is itself a vector space.

\medskip

\textbf{Question.} Do polynomials of degree exactly $n$ form a subspace?

\medskip

No. If
\[
a_0+a_1x+\cdots+a_nx^n
\]
has degree $n$, then subtracting $a_nx^n$ gives
\[
a_0+a_1x+\cdots+a_{n-1}x^{n-1},
\]
which does not have degree $n$.

\medskip

\textbf{Question.} What is a basis of $\mathbb{R}_{\leq n}[x]$?

\[
\beta=\{1,x,x^2,\ldots,x^n\}.
\]

\newpage

\textbf{Example.} Do upper triangular $n\times n$ matrices form a subspace of the vector space of all $n\times n$ matrices?

\medskip

Yes. For example,
\[
\begin{bmatrix}1&2\\0&3\end{bmatrix}
+5\begin{bmatrix}-1&1\\0&2\end{bmatrix}
\]
is upper triangular; addition and scalar multiplication preserve the zero entries below the diagonal.

\bigskip

\textbf{Example.} A basis of $M_{m\times n}(\mathbb{R})$, the vector space of $m\times n$ matrices, consists of the matrices having exactly one nonzero entry, equal to $1$.

For instance,
\[
\begin{bmatrix}2&1&3\\5&1&7\end{bmatrix}
=2\begin{bmatrix}1&0&0\\0&0&0\end{bmatrix}
+\begin{bmatrix}0&1&0\\0&0&0\end{bmatrix}
+\cdots.
\]
Therefore,
\[
\dim M_{m\times n}(\mathbb{R})=mn.
\]

\bigskip

\textbf{Definition.} The \textbf{dimension} of a vector space is the number of vectors in a basis.

\medskip

\textbf{Examples.}
\[
\dim\mathbb{R}_{\leq n}[x]=n+1,
\]
because its elements have the form $a_0+a_1x+\cdots+a_nx^n$.

\[
\dim\mathbb{R}[x]=\infty,
\]
because $\{1,x,x^2,\ldots,x^n,\ldots\}$ is a basis.

\newpage

\textbf{Examples.}
\[
\dim M_{n\times n}(\mathbb{R})=n^2.
\]

The vector space of $n\times n$ upper triangular matrices has dimension
\[
n+\frac{n^2-n}{2}=\frac{n^2+n}{2}.
\]

\bigskip

\textbf{Definition.} A map $T:V\to W$, where $V,W$ are vector spaces, is \textbf{linear} if
\[
T(c_1v_1+c_2v_2)=c_1T(v_1)+c_2T(v_2)
\]
for every $v_1,v_2\in V$ and $c_1,c_2\in\mathbb{R}$.

\medskip

\textbf{Example.} Define
\[
T:M_{n\times n}(\mathbb{R})\to M_{n\times n}(\mathbb{R}),
\qquad T(A)=A^T.
\]
Then
\[
T(c_1A_1+c_2A_2)=(c_1A_1+c_2A_2)^T
=c_1A_1^T+c_2A_2^T
=c_1T(A_1)+c_2T(A_2),
\]
so $T$ is linear.

\newpage

\textbf{Example.} Define
\[
T:\mathbb{R}[x]\to\mathbb{R}[x],
\qquad T(f)=f'.
\]
Then
\[
T(c_1f_1+c_2f_2)=(c_1f_1+c_2f_2)'
=c_1f_1'+c_2f_2'
=c_1T(f_1)+c_2T(f_2).
\]
Thus $T$ is linear.

\medskip

\textbf{Example.} For this derivative map,
\[
\ker T=\{f\in\mathbb{R}[x]:f'=0=T(f)\},
\]
which is the set of constant polynomials.

\medskip

\textbf{Example.} Define
\[
T:\mathbb{R}[x]\to\mathbb{R}[x],
\qquad
T(f)=\int_c^x f(t)\,dt.
\]
Since
\[
\int_c^x(c_1f_1+c_2f_2)\,dt
=c_1\int_c^x f_1\,dt+c_2\int_c^x f_2\,dt,
\]
$T$ is linear.

\newpage

\textbf{Definition.} A vector space with a finite basis is \textbf{finite-dimensional}.

\bigskip

\textbf{Theorem.} Let
\[
\beta=\{v_1,\ldots,v_n\}
\]
be a basis of a vector space $V$. Let $T:V\to\mathbb{R}^n$ be the coordinate map
\[
T(x)=[x]_\beta.
\]
That is,
\[
T(c_1v_1+\cdots+c_nv_n)=
\begin{bmatrix}c_1\\ \vdots\\ c_n\end{bmatrix}.
\]
Then $T$ is linear and invertible.

\bigskip

\textbf{Definition.} A linear invertible map is a \textbf{linear isomorphism}.

\medskip

If $T:V\to W$ is a linear isomorphism, then $V$ and $W$ are \textbf{isomorphic}.

\medskip

\textbf{Example.} The map
\[
a_0+a_1x+\cdots+a_nx^n
\longmapsto
\begin{bmatrix}a_0\\ \vdots\\ a_n\end{bmatrix}
\]
is a linear isomorphism from $\mathbb{R}_{\leq n}[x]$ to $\mathbb{R}^{n+1}$. The basis is
\[
\beta=\{1,x,\ldots,x^n\}.
\]

\newpage

\textbf{Key point.} In isomorphic vector spaces, linear algebra works in the same way.

\medskip

\textbf{Theorem.} If $T:V\to W$ is linear and invertible, then
\[
T^{-1}:W\to V
\]
is linear.

\medskip

\textbf{Proof.} Suppose, for a contradiction, that
\[
T^{-1}(c_1w_1+c_2w_2)\ne c_1T^{-1}(w_1)+c_2T^{-1}(w_2).
\]
Since $T$ is one-to-one, applying $T$ preserves the inequality. But
\begin{align*}
T\bigl(T^{-1}(c_1w_1+c_2w_2)\bigr)
&=c_1w_1+c_2w_2\\
&=c_1T\bigl(T^{-1}(w_1)\bigr)+c_2T\bigl(T^{-1}(w_2)\bigr)\\
&=T\bigl(c_1T^{-1}(w_1)+c_2T^{-1}(w_2)\bigr),
\end{align*}
which is a contradiction.

\newpage

\textbf{Definition.} Let $S$ be the space of sequences $(x_k)_{k\geq0}$ with $x_k\in\mathbb{R}$. Define
\[
\operatorname{Fib}=
\left\{(x_k)\in S:x_{k+2}=x_{k+1}+x_k\right\}.
\]

\medskip

For example, if $(x_k)\in\operatorname{Fib}$, then
\[
2(x_k)=2(x_k),
\qquad
2x_{k+2}=2x_{k+1}+2x_k,
\]
so scaling preserves the recurrence. Likewise, addition preserves the recurrence. Thus $\operatorname{Fib}$ is a subspace of $S$.

\medskip

Define
\[
T:\operatorname{Fib}\to\mathbb{R}^2,
\qquad
(x_k)\longmapsto
\begin{bmatrix}x_0\\x_1\end{bmatrix}.
\]
This map is linear and invertible: the recurrence determines every later term, for example,
\[
x_2=x_0+x_1,
\qquad
x_3=x_1+x_2.
\]
Therefore $\operatorname{Fib}$ is linearly isomorphic to $\mathbb{R}^2$, and a basis can be obtained from
\[
\{T^{-1}(e_1),T^{-1}(e_2)\}.
\]

\newpage

To find vectors in $\operatorname{Fib}$, search for sequences of the form $r^k$, where $r\ne0$. The recurrence gives
\[
r^{k+2}=r^{k+1}+r^k,
\]
so
\[
r^2-r-1=0,
\qquad
r_{\pm}=\frac{1\pm\sqrt5}{2}.
\]
Thus
\[
(r_+^k),\qquad(r_-^k)
\]
belong to $\operatorname{Fib}$ and are independent, so they form a basis of $\operatorname{Fib}$.

\medskip

Every Fibonacci sequence has the form
\[
f_k=C_+\left(\frac{1+\sqrt5}{2}\right)^k
+C_-\left(\frac{1-\sqrt5}{2}\right)^k.
\]
For the sequence with $f_0=0$ and $f_1=1$,
\[
0=C_++C_-,
\qquad
1=C_+r_++C_-r_-.
\]

\end{document}
